\title{Byte Latent Transformer: Patches Scale Better Than Tokens}

\begin{abstract}
We present the Byte Latent Transformer ({\textsc BLT}), a novel large language model architecture operating at the byte level that, for the first time, achieves parity in performance with tokenization-based LLMs at scale, while also offering considerable gains in inference speed and robustness. {\textsc BLT} operates by encoding input bytes into variable-length patches, which function as its core computational units. The segmentation of patches is guided by estimating the entropy of the upcoming byte, thereby dedicating additional model resources to segments exhibiting greater complexity. We conduct the first flop-constrained scaling study on byte-level models extending up to 8B parameters and 4T bytes of training data. Our experiments confirm that it is practical to scale models on raw byte sequences without relying on a predefined vocabulary. Dynamically leveraging extended patches in low-entropy regions leads to improvements in both training and inference efficiency, as well as enhanced performance on tasks demanding reasoning and generalization to rare cases. In summary, {\textsc BLT} attains a superior scaling relationship compared to tokenization-based LLMs at fixed inference costs by jointly expanding both patch size and model capacity.
\end{abstract}

\section{Introduction}
\label{section:intro}

We present the Byte Latent Transformer~(\textbf{BLT}), a novel architecture that eliminates the need for tokenizers by directly leveraging raw byte inputs. BLT, for the first time, achieves parity in performance with conventional tokenization-based models at large scales while also delivering notable gains in both computational efficiency and resilience to noise (\S\ref{sec:robustness}). 
Currently, most large language models (llms) are trained in an almost fully end-to-end fashion, with the exception of the tokenization step—a heuristic process that segments byte sequences into a predefined set of tokens. This step introduces inductive biases during string compression, which can result in problems such as sensitivity to specific domains or modalities~\citep{dagangetting}, vulnerability to noisy inputs~(\S\ref{sec:robustness}), limited orthographic awareness~\citep{edman2024cute}, and unequal multilingual representation~\citep{liang2023xlm,petrov2024language,limisiewicz2024myte}.

Historically, tokenization has played a crucial role since training large language models on raw bytes is prohibitively expensive at scale owing to the increased sequence lengths~\citep{xue2022byt5}. To address this, previous research has utilized more efficient self-attention mechanisms~\citep{el2020characterbert, clark2022canine} or explored architectures that bypass attention altogether~\citep{wang2024mambabyte}~(\S\ref{section:related_work}). Nonetheless, these strategies predominantly benefit the training of \textit{smaller models}. When scaling up, the primary computational expense in a Transformer is attributed to the extensive feed-forward network layers applied to every byte, while the attention mechanism itself represents a minor portion of the overall cost.

To enable more efficient use of computation, we introduce a dynamic and learnable strategy that clusters bytes into \textit{patches}~(\S\ref{section:patching}), alongside a novel architecture that fuses both byte-level and patch-level information. In contrast to conventional tokenization, {\textsc BLT} does not rely on a static vocabulary for patches. Instead, any combination of bytes can be transformed into latent patch representations through compact, learnable encoder and decoder components. Our findings demonstrate that this approach leads to a \textit{greater} efficiency in the distribution of compute resources compared to models based on traditional tokenization.

Tokenization-based LLMs distribute computation uniformly across all tokens, regardless of their predictive complexity. This uniform allocation compromises efficiency because the tokenization process relies on compression strategies that do not necessarily align with the varying difficulty of prediction tasks. Our architecture is fundamentally built on the principle that computational resources should be applied adaptively, concentrating effort where it is most beneficial. For instance, predicting the final character in most words is a relatively low-entropy, straightforward operation and does not require the capacity of a large transformer; in contrast, generating the initial word of a sentence often involves more complex, high-entropy decisions. Reflecting this insight, {\textsc BLT} adopts an architecture~(\S\ref{section:architecture}) featuring three transformer components: two lightweight local models at the byte level, and a more powerful global \textit{latent transformer}~(\autoref{fig:arch_high_level}). To enable adaptive computation, {\textsc BLT} organizes bytes into patches by assessing the entropy of the next-byte predictions, thereby forming contextualized segments with approximately consistent information content.

We conduct the first study examining flop-controlled scaling for byte-level architectures, scaling up to 8B parameters and processing 4T bytes in training. Our results demonstrate the feasibility of large-scale, end-to-end byte-level model training without reliance on predefined vocabulary tokenization. In terms of training efficiency, {\textsc BLT} achieves parity with Llama 3 under identical flop constraints,\footnote{Model computational cost is measured by tallying total Floating Point OPerations (flops).} while reducing inference flops by as much as 50\% (\S\ref{section:scaling}).

Moreover, using raw byte input substantially improves performance on rare data distributions, where traditional approaches struggle. {\textsc BLT} models exhibit greater robustness to input noise and excel in fine-grained character-level processing, as evidenced by results on orthographic, phonological, and low-resource machine translation benchmarks (\S\ref{sec:robustness}).

Additionally, the {\textsc BLT} framework enables concurrent scaling of both model capacity and patch size, all under a consistent inference flop limit. Increasing patch size typically reduces overall computation, permitting resource savings to be directed toward enlarging the global latent transformer, since it executes less frequently. Our experiments, detailed in (\autoref{fig:fixed_inference_scaling}), reveal markedly improved scaling properties compared to conventional tokenization-based models.

To conclude, the primary contributions of this work are as follows: 1) We present {\textsc BLT}, a byte latent llm framework capable of dynamically distributing computational resources to enhance flop efficiency, 2) Our results indicate that we match the training flop requirements of Llama 3 at scales up to 8B, while allowing for strategic trade-offs where modest declines in evaluation performance yield up to 50\% improvements in flop efficiency, 3) {\textsc BLT} models enable a novel scaling regime for llms, permitting growth in model size without exceeding a constant inference compute budget, and 4) We provide evidence that {\textsc BLT} models exhibit superior resilience to noisy inputs and exhibit an understanding of sub-word input features not typically captured by token-level llms. The codebase for training and inference using {\textsc BLT} is openly available at~\url{https://github.com/facebookresearch/blt}.

\section{Patching: From Individual Bytes to Groups of Bytes}
\label{section:patching}

Dividing byte sequences into \textit{patches} enables {\textsc BLT} to flexibly assign computational resources in response to varying contextual demands. Figure~\ref{fig:patching_types} illustrates multiple strategies for segmenting bytes into patches. More formally, a patching function $f_p$ operates on a byte sequence $\pmb{x}=\{x_i,|i=1,\ldots n\}$ of length $n$, partitioning it into $m < n$ patches $\pmb{p}=\{p_j|j=1,\ldots,m\}$, where each $x_i$ is associated with either 0 or 1—1 marking the beginning of a new patch. For both token-based and patch-based approaches, the main Transformer’s computational expenditure is driven by the number of main processing steps, corresponding in {\textsc BLT} to the count of patches that the patching function yields when encoding data. Thus, the typical or mean \textit{patch size} (i.e., average bytes per patch) becomes the central determinant of computational requirements in both model training and inference under any given patching method~(\S\ref{section:flops}).

In what follows, we describe three distinct patching functions: fixed-size patching, where patches consist of a constant number of bytes~(\S\ref{section:static-patch}); whitespace-based patching, which segments at spaces or similar delimiters~(\S\ref{section:white-patch}); and entropy-based dynamic patching, where patch boundaries are guided by local entropy as estimated by a lightweight byte language model~(\S\ref{section:dyn-patch}). Lastly, we examine the idea of incremental patching and outline the distinctions between traditional tokenization and patch-based segmentation~(\S\ref{section:bpe-patch}).

\subsection{Strided Patching Every K Bytes}
\label{section:static-patch}

A common approach to aggregating bytes involves organizing them into fixed-size patches of length $k$, as illustrated in MegaByte~\citep{yu2023megabyte}. Using a constant stride simplifies both training and inference procedures, and allows for direct adjustment of average patch size—thereby facilitating the management of computational cost. Nonetheless, this type of patching entails several notable drawbacks. Notably, computational resources are not adaptively assigned where most needed; a transformer step $j$ might be wasted if only whitespace in source code is being modeled, or conversely, insufficient compute might be dedicated to bytes containing highly informative content, such as mathematics. Furthermore, this strategy produces inconsistent, non-contextual segmentations of identical byte sequences—for instance, different splits of the same word across occurrences.

\subsection{Space Patching}
\label{section:white-patch}

According to \citet{slagle2024spacebyte}, an improved method is presented as an alternative to strided patching, which generates new patches immediately after encountering any space-like byte\footnote{
    Space-like bytes refer to any byte not corresponding to a Latin character, digit, or utf-8 continuation byte. Furthermore, at least one non space-like byte must be present in each patch. }—positions that frequently demarcate linguistic boundaries in various languages. In this approach, called Space patching, an additional latent transformer operation (i.e., increased computational resources) is dedicated to every individual word, standardizing the patching of words across input sequences and targeting computational effort towards challenging prediction points that often succeed space-like bytes. For instance, generating the initial byte of the response to ``Who composed the Magic Flute?\uline{\hspace{.6em}}'' poses a considerably more complex task compared to producing subsequent bytes in ``Mozart'' after knowing ``M''; the first letter drastically narrows plausible continuations, simplifying the remaining prediction. Despite these advantages, space patching exhibits limitations, notably its inability to flexibly adjust patch sizes and its suboptimal handling of certain languages and application domains. Therefore, in the following, we present a novel patching strategy based on the observation that the initial bytes of words tend to be the most challenging to predict, offering a principled method for managing patch sizes.

\subsection{Entropy Patching: Using Next-Byte Entropies from a Small Byte LM}
\label{section:dyn-patch}

Instead of depending on heuristic rules like whitespace for boundary identification, we adopt a data-driven method for detecting areas with high uncertainty in next-byte predictions. We propose \textit{entropy patching}, a technique that leverages entropy measurements to determine patch boundaries.

We train a small byte-level auto-regressive language model on the training data for {\textsc BLT} and compute next byte entropies under the LM distribution $p_{e}$ over the byte vocabulary $\mathcal{V}$:

\begin{equation}
H(x_i) = \sum_{v \in \mathcal{V}} p_{e}(x_i=v|\pmb{x}_{<i}) \log p_{e}(x_i=v|\pmb{x}_{<i})
\end{equation}

We utilize two approaches to detect patch boundaries from the entropy values $H(x_i)$. The initial method selects points where the entropy exceeds a fixed global threshold, as demonstrated in~\autoref{fig:patching}. The alternative technique flags points that show a substantial increase compared to the preceding entropy value. This latter method effectively captures locations that disrupt the general trend of decreasing entropy within a given patch.

\begin{align*}
    \text{Global Constraint}\hspace{1cm}H(x_t)& > \theta_g\\
    \text{Approx. Monotonic Constraint}\hspace{1cm}H(x_t)& - H(x_{t-1}) > \theta_r
\end{align*}

Patch boundaries are determined through a streamlined preprocessing phase carried out while loading data. This contrasts with~\citet{nawrot-etal-2023-efficient}, where a classifier is employed to predict entropy-based patch boundaries. Within our experiments~(\S\ref{section:experiments}), we perform a comparison between these two strategies for differentiating low and high entropy bytes.

\subsection{The Byte-Pair Encoding (BPE) Tokenizer and Incremental Patching}
\label{section:incremental-patching}
\label{section:bpe-patch}

Numerous contemporary llms, such as our baseline Llama 3, rely on subword tokenizers like bpe~\citep{gage1994new, sennrich-etal-2016-neural}. In this work, we use the term ``tokens'' to mean byte groupings selected from a \textit{finite} vocabulary defined before training, whereas ``patches'' denote dynamically created sequences that are not limited by a preset vocabulary. Importantly, a key distinction is that, when using tokens, the model cannot directly utilize the original byte-level information.

An important advancement offered by {\textsc BLT} relative to tokenization-based models is its reconsideration of the balance between vocabulary size and computational resources. In conventional llms, expanding the vocabulary leads to longer tokens on average, which results in a reduced number of model steps but simultaneously enlarges the output dimensionality required for the model’s final projection layer. This compromise generally constrains tokenization-based methods from enabling substantial variability in both token length and inference efficiency. As an illustration, Llama 3 raises the average token length from 3.7 to 4.4 bytes, but this comes at the expense of making its embedding table four times larger than that of Llama 2.

During generation, {\textsc BLT} must determine at each step of the byte sequence whether it has reached a patch boundary, which in turn decides if additional computation is required through the Latent Transformer. Importantly, this determination should be made without relying on any future portions of the sequence that have not yet been produced. As such, the patching process cannot depend on access to subsequent bytes when establishing the segmentation of the byte sequence. More formally, a patching method $f_p$ is said to implement incremental patching if it holds that:
$$f_p(\pmb{x}_{<i}) = f_p(\pmb{x})_{<i}$$
for all positions $i$. In contrast, bpe does not fulfill the criteria for incremental patching, because an identical prefix can lead to different tokenizations based on the upcoming sequence, and thus, it fails to meet the property outlined above\footnote{It is possible to convert bpe into an incremental patching method by introducing a special delimiter token to signal patch boundaries, though this would increase the overall length of the byte sequence.}.

\section{{\textsc BLT} Architecture}
\label{section:architecture}
The {\textsc BLT} system comprises a central, large-scale autoregressive language model designed to process patch-level representations, supplemented by two lightweight local modules responsible for transforming sequences of bytes to patches and reconstructing byte sequences from patch representations~(\autoref{fig:arch_high_level}).

\subsection{Latent Global Transformer Model}

\textit{The Latent Global Transformer} refers to an autoregressive transformer model, denoted as $\mathcal{G}$, comprising $l_{\mathcal{G}}$ layers. This model takes a sequence of latent input patch embeddings, $p_j$, and produces a corresponding sequence of output patch embeddings, $o_j$. For clarity, subscripts $j$ and $i$ are consistently used to represent patches and bytes, respectively, throughout the paper. The global transformer employs a block-causal attention mask~\citep{dubey2024llama}, which confines attention to patches up to and including the current one within any given document. The majority of computational effort, both during pre-training and inference, is attributed to this model. Consequently, strategically selecting when to apply the global model enables flexible management of computational resources, allowing compute allocation to vary according to the complexity of specific input or output regions.

\subsection{Local Encoder}
\label{section:local}

\textit{The Local Encoder Model}, symbolized as $\mathcal{E}$, is designed as a compact transformer architecture, employing $l_{\mathcal{E}}$ layers where $l_{\mathcal{E}} << l_{\mathcal{G}}$. Its primary function is to rapidly convert a sequence of input bytes $b_i$ into rich patch-level representations $p_j$. A notable innovation compared to standard transformer architectures is the incorporation of a cross-attention module immediately after each transformer block. This module serves to aggregate the byte-level embeddings into patch-level features~(\autoref{fig:crossattn}).

Initially, each input byte $b_i$ is embedded via a $\mathbb{R}^{256 \times h_{\mathcal{E}}}$ embedding matrix, resulting in vectors $x_i$. These base embeddings can be further enhanced by integrating hash-based embedding features as detailed in~(\S\ref{sec:hash-emb}). The representations are then passed through a stack that alternates between transformer and cross-attention blocks~(\S\ref{sec:enc-cross-attn}), systematically refining them into patch embeddings $p_i$. These patch representations subsequently serve as input to the global transformer module $\mathcal{G}$.

Each transformer layer makes use of a \textit{local block causal} attention mask: each position in the byte sequence can attend to a restricted window, specifically the preceding $w_{\mathcal{E}}$ bytes. This local attention can span patch boundaries as they change dynamically but does not extend beyond the confines of the current document. Further architectural components, such as embedding strategies and the cross-attention mechanism, are elaborated in the following subsections.

\subsubsection{Encoder Hash n-gram Embeddings}
\label{sec:ngram-emb}
\label{sec:hash-emb}
An essential factor in producing effective and informative representations at every position $i$ is to include context from the prior bytes. In {\textsc BLT}, this is accomplished by jointly encoding the byte $b_i$ on its own and within the context of a byte n-gram. At each index $i$, we initially generate byte-grams

\begin{align}
    g_{i,n}=\{b_{i-n + 1},\ldots, b_i\}
\end{align}

for each byte index $i$ and for all $n$ between three and eight.\footnote{
Byte-grams of length $n$ or greater are omitted when $i<n$. }

Next, we define hash $n$-gram embeddings that use a hash function to project every byte $n$-gram to a slot in a fixed-size embedding matrix $E_{n}^{hash}$ corresponding to each $n$ in $\{3,4,5,6,7,8\}$~\citep{bai2010learning}. This hash-derived embedding is then summed with the byte's embedding, after which the sum is normalized and utilized as input to the local encoder. The final augmented representation is computed as:

\begin{align}
e_i &= x_i + \sum_{n = 3,...,8}E_{n}^{hash}(\text{Hash}(g_{i,n})) \\
\text{where, Hash}(g_{i,n}) &= \text{RollPolyHash}(g_{i,n}) \% |E_{n}^{hash}|
\end{align}

The value $e_i$ is divided by the total number of $n$-gram lengths plus one for normalization, and the RollPolyHash technique described in Appendix~\ref{appendix:rollpolyhash} is utilized. Section~\ref{section:ablations} presents an ablation study assessing how variations in $n$-gram hash embeddings, specifically altering $n$ and the embedding table size, affect flop-controlled scaling law patterns. Beyond hash-based $n$-gram embeddings, we also investigated frequency-based $n$-gram embeddings; further information regarding these experiments can be found in Appendix~\ref{appendix:ngrams}.

\subsubsection{Encoder Multi-Headed Cross-Attention}
\label{sec:enc-cross-attn}
Our implementation is largely inspired by the input cross-attention mechanism in the Perceiver architecture~\citep{jaegle2021perceiver}, though it introduces a key modification: instead of mapping a fixed set of latent representations, we map latent variables to the dynamic patch representations (\autoref{fig:crossattn}), restricting each to attend solely to the bytes contained within their respective patch. In this module, each patch $p_j$ is associated with a query vector, which is derived by aggregating the byte-level representations for that specific patch via pooling, and subsequently passing this pooled output through a linear projection layer, $\mathcal{E}_{C} \in \mathbb{R}^{h_{\mathcal{E}} \times (h_{\mathcal{E}} \times U_{\mathcal{E}})}$. Here, $U_{\mathcal{E}}$ is the number of distinct attention heads in the encoder's cross-attention mechanism. More formally, let $f_{\text{bytes}}(p_j)$ indicate the byte sequence for patch $p_j$. The resulting computations are then given by

\begin{align}
P_{0,j} &= \mathcal{E}_{C}(f_\text{bytes}((p_j)), f~ \text{is a pooling function} \\
P_l &= P_{l-1} + W_o\left(\text{softmax}\left(\frac{QK^T}{\sqrt{d_k}}\right)V\right) \\
\text{where } Q_j &= W_q(P_{l-1,j}), K_i = W_k(h_{l-1,i}), V_i = W_v(h_{l-1,i}) \\
h_l &= \text{Encoder-Transformer-Layer}_l(h_{l-1})
\end{align}

In this framework, $P \in \mathbb{R}^{n_p \times h_{\mathcal{G}}}$ denotes the set of $n_p$ patch-level representations that serve as inputs to the global model. These are initially formed by aggregating the byte embeddings $e_i$ associated with each patch $p_j$. The matrices $W_q$, $W_k$, $W_v$, and $W_o$ are linear projections for queries, keys, values, and outputs, respectively. The keys and values are generated by projecting the byte representations $h_i$ from the preceding layer (or $e_i$ at the initial layer). To enforce locality, a tailored masking technique is implemented, such that each query $Q_j$ is restricted to attend solely to the keys and values derived from the bytes within its corresponding patch $j$. Since multi-headed attention operates over $Q, K,$ and $V$, and patch embeddings generally have a higher dimensionality ($h_{\mathcal{G}}$) than $h_{\mathcal{E}}$, we represent $P_l$ across multiple attention heads, each of size $h_{\mathcal{E}}$ in the cross-attention computation, and subsequently concatenate these to achieve the required $h_{\mathcal{G}}$ dimensionality. Furthermore, pre-LayerNorm normalization is applied to the queries, keys, and values, and this cross-attention layer does not incorporate positional embeddings. A residual connection is finally applied around the cross-attention component to stabilize training.

\subsection{Local Decoder} 

Analogous to the local encoder, the local decoder $\mathcal{D}$ employs a lightweight transformer architecture comprising $l_{\mathcal{D}} << l_{\mathcal{G}}$ layers. Its role is to transform a sequence of global patch embeddings, $o_j$, back into a sequence of raw bytes, $y_i$. The local decoder generates these raw byte predictions conditioned on the sequence of bytes decoded up to that point, utilizing the hidden representations output by the local encoder for the entire byte sequence. This decoding process involves $l_{\mathcal{D}}$ successive layers, each alternately performing cross-attention and transformer operations. In each layer, cross-attention is executed prior to the transformer block, enabling the construction of byte-level representations based on the input patch representations; subsequently, the local transformer processes this derived byte sequence.

\subsubsection{Decoder Multi-headed Cross-Attention} 

In the decoder cross-attention, the roles of the queries and key/values are interchanged i.e. the byte-representations are now the queries, and the patch representations are now the key/values. The initial byte-representations for the cross-attention are initialized as the byte embeddings from the last encoder layer i.e. $h_{l_{\mathcal{E}}}$. The subsequent byte-representations for layer $l$, $d_{l,i}$ are computed as:

\begin{align}
D_0 &= h_{l_{\mathcal{E}}} \\
B_l &= D_{l-1} + W_o\left(\text{softmax}\left(\frac{QK^T}{\sqrt{d_k}}\right)V\right), \\
  &\text{where } Q_i = W_q(d_{l-1,i}), K_i = W_k(\mathcal{D}_C(o_j)), V_i = W_v(\mathcal{D}_C(o_j)) \\
  D_l &= \text{Decoder-Transformer-layer}_l(B_l)
\end{align}

Here, as before, $W_k$ and $W_v$ denote the key and value projection matrices, which perform a linear transformation followed by a split operation $\mathcal{D}_C$ on the final patch representations $o_j$ coming from the global model. The query projection matrix $W_q$ operates over the byte-level representations $d_{l-1}$ obtained from the prior decoder transformer layer (or over $h_{l_{\mathcal{E}}}$ when initializing at the first layer). The output projection is done by $W_o$, so that $B \in \mathbb{R}^{h_{\mathcal{D}} \times n_b}$, where $n_b$ designates the count of output bytes. The subsequent layer of decoder representations $D_l$ are then computed by applying a decoder transformer layer to the results of the cross-attention module $B$. As with the cross-attention in the local encoder, the model employs multi-head attention, includes pre LayerNorms, omits positional embeddings, and incorporates a residual connection that encircles the cross-attention operation.

\section{Experimental Setup}
\label{section:experiments}
We structure controlled experimental conditions to benchmark {\textsc BLT} against models relying on tokenization, specifically ensuring that {\textsc BLT} does not benefit from potential access to extended sequence context.

\subsection{Pre-training Datasets} In this work, all tested model configurations undergo pre-training on two distinct datasets: (1) the Llama 2 dataset~\citep{touvron2023llama}, which contains 2 trillion tokens sourced from a wide array of publicly accessible data that has been rigorously cleaned and filtered to enhance data quality; and (2) {\textsc BLT}-1T, a newly curated dataset composed of 1 trillion tokens drawn from multiple public data sources and partially including segments from the pre-training set introduced by Datacomp-LM~\citep{li2024datacomp}. The Llama 2 dataset primarily serves for scaling law investigations, specifically to ascertain the optimal token counts for pre-training according to the findings of~\cite{dubey2024llama}, thus informing architectural design decisions for {\textsc BLT}. Conversely, {\textsc BLT}-1T is leveraged for end-to-end pre-training experiments to facilitate direct comparisons against Llama 3 on a range of downstream benchmarks. Importantly, both datasets are devoid of any information derived from Meta products or services. Additionally, for baseline assessments involving tokenizer-based models, we employ the Llama 3 tokenizer, which features a vocabulary comprising 128K tokens; this tokenizer consistently yielded stronger baselines compared to the Llama 2 tokenizer throughout our study.

\subsection{Entropy Model} In our study, the entropy model is implemented as a byte-level language model that is trained on the identical data distribution as the comprehensive {\textsc BLT} model. By default, the architecture consists of a transformer with 100M parameters, featuring 14 layers, a hidden size of 512, and utilizes sliding window attention over a context of 512 bytes. All other hyperparameters match those of both our local and global transformer variants. Variations in model size, receptive field, and architecture were investigated, as outlined in \autoref{section:ablations}. Notably, with sufficiently limited receptive fields, it is possible to represent the trained entropy model using a compact lookup table.

\subsection{Entropy Threshold and Equalizing Context Length} In entropy-based patching methods, we determine a threshold intended to yield a specific average \textit{patch size} on the pretraining dataset mix. In contrast to tokenization, {\textsc BLT} allows for arbitrary selection of \textit{patch size}, which directly affects the model’s context length. To prevent bias that might benefit models with larger patch sizes through longer contexts, we fix the expected number of bytes per batch across all configurations. Consequently, when increasing patch size, we proportionally decrease sequence length to maintain equivalence. For Llama 2 data, models use an 8k byte context window, whereas on the {\textsc BLT}-1T dataset, we extend the average context to 16k bytes, all while preserving a consistent batch size of 16M bytes on average.

Although the average batch size remains unchanged, dynamic patching techniques produce varying byte-to-patch ratios during batch loading. In our approach to {\textsc BLT} training, we optimize efficiency by assembling batches to contain an equal number of patches, which removes the need for padding within the costly latent transformer component. To manage memory usage and prevent spikes resulting from sequences containing exceptionally large patches, we pad or, when necessary, truncate input byte sequences to 12k and 24k bytes for the Llama 2 and {\textsc BLT}-1T datasets, respectively.

\subsection{Entropy Model Context}
\label{section:entropy-context}
Our empirical observations reveal that entropy patching tends to produce increasingly larger patches within highly structured data, such as multiple choice questions (see example in \autoref{fig:mmlu_patching}), where content is frequently repetitive. These differences can be attributed to reduced entropy associated with recurring material present in the entropy model context. Consequently, for the large-scale execution of {\textsc BLT}-Entropy using a patch size of 4.5, we opted to reset the entropy context using line breaks and imposed an approximate monotonicity constraint. This approach is more resilient to the effects of “entropy drift” that arise from fluctuations in context length. Notably, this adjustment pertains solely to our entropy calculations; the methodology for determining the entropy threshold remains unchanged.

\subsection{FLOPs Estimation} 
\label{section:flops}

Our methodology for calculating transformer FLOPs adheres closely to the formulae provided by Chinchilla~\citep{hoffmann2022training}, incorporating the FLOPs attributed to the feed-forward modules, the qkvo projections in self-attention, as well as the operations for attention computation and output projection. Unlike their setup, however, we model the input embedding layer as a computationally inexpensive lookup rather than as a dense matrix multiplication, effectively rendering its FLOP cost negligible. Consistent with prior literature, we approximate that the backward pass incurs a computational load twice that of the forward pass.

To compute flops \textit{per byte} for {\textsc BLT} models, we add up the flops for the local encoder transformer, the global latent transformer, and the local decoder transformer, together with the cross attention blocks in the encoder and the decoder:

\begin{align}
    \text{FL}_{\text{{\textsc BLT}}} &= \text{Transf. FL}(h_{\mathcal{G}}, l_{\mathcal{G}}, m=n_{ctx}/n_p, V=0) / n_p\\
   &\quad +\text{Transf. FL}(h_{\mathcal{E}}, l_{\mathcal{E}}, m=w_{\mathcal{E}}, V=0) \\
   &\quad + \text{Transf. FL}(h_{\mathcal{D}}, l_{\mathcal{D}}, m=w_{\mathcal{D}}, V=256) \\ 
    &\quad + \text{Cross Attn. FL}(h_{\mathcal{E}}, l_{\mathcal{E}}, m=n_p, r=n_p/k) \cdot \frac{k}{n_p} \\ 
   &\quad + \text{Cross Attn. FL}(h_{\mathcal{D}}, l_{\mathcal{D}}, m=k, r=k/n_p) 
\end{align}

Here, $n_{ctx}$ denotes the input sequence length measured in bytes, $n_p$ represents the patch size, $r$ specifies the proportion of queries relative to keys/values, and $k$ indicates the ratio between patch dimension and byte dimension—effectively the number of local model partitions that are joined to yield the overall model representation ($k=2$ in \autoref{fig:crossattn}). The term $V$ signifies the size of the vocabulary in the output projection layer, a component unique to the local decoder. The context length used in attention, denoted by $m$, differs according to whether the module processes byte sequences or patch sequences. As such, we revise the attention FLOPs formula for each module. Detailed formulations for calculating Transformer-FLOPs as well as Cross-Attention FLOPs appear in Appendix~\ref{section:flops-appendix}.

\subsection{Bits-Per-Byte Estimation} Perplexity only makes sense in the context of a fixed tokenizer as it is a measure of the uncertainty for each token. When comparing byte and token-level models, following previous work~\citep{xue2022byt5, yu2023megabyte, wang2024mambabyte}, we instead report Bits-Per-Byte (BPB), a tokenizer independent version of perplexity. Specifically:

\begin{align}
    \text{BPB}(x)&= \frac{\mathcal{L}_{CE}(\pmb{x})}{\ln(2)\cdot n_{\text{bytes}}}
\end{align}

In this formulation, the cross-entropy loss accumulated over the data $\pmb{x}$, representing its uncertainty, is divided by both the total count of bytes in $\pmb{x}$ and a normalization factor.

\subsection{Transformer Architecture Hyperparameters} 

Across all transformer modules in {\textsc BLT}—encompassing both local and global networks—we adopt an architectural approach akin to Llama~3~\citep{dubey2024llama}. Specifically, our feed-forward components employ the SwiGLU activation~\citep{swiglu}, while the self-attention layers utilize rotary positional encodings (RoPE)~\citep{RoPe} with $\theta=500000$ as described by~\citep{xiong2024effective}. Layer normalization is conducted with RMSNorm~\citep{RMSNorm}. To optimize memory and computation, Flash Attention~\citep{dao2022flashattention} is leveraged for all self-attention operations governed by pre-defined masks, such as \textit{block causal} or \textit{fixed-window block causal}, with fixed-width masks using a window span of 512 tokens. For the cross-attention modules, which require dynamically determined patch-sensitive masks, we incorporate Flex Attention\footnote{\url{https://pytorch.org/blog/flexattention}} to generate efficient, fused kernels, thereby considerably accelerating the training process.

\subsection{{\textsc BLT}-Specific Hyperparameters}  
To evaluate the performance of {\textsc BLT} models, we design experiments focusing on both scaling patterns and downstream tasks, exploring model sizes of 400M, 1B, 2B, 4B, and 8B parameters. Architectural hyperparameter details for these models can be found in Appendix~Table~\ref{tab:arch}.  
For the initialization of the queries in the local encoder's initial cross-attention layer, we apply max-pooling. The hashing setup utilizes $500,000$ hashes generated by a single hash function, and we consider n-grams in the range of sizes from 3 to 8 across all {\textsc BLT} models. The learning rate is consistently fixed at $4e-4$ for all configurations. The rationale for adopting this rate for both token-based and {\textsc BLT} architectures originates from a hyperparameter search over the interval $1e-3$ to $1e-4$, where results from the 400M and 1B scales confirmed this value as optimal. In our scaling analysis utilizing Llama-2 data, we determine training batch sizes according to recommendations in~\cite{dubey2024llama}, or set equivalently in terms of data volume (bytes). For model training, AdamW~\citep{adamw} is employed as the optimizer, configuring $\beta_1$ at 0.9, $\beta_2$ at 0.95, and $\epsilon=10^{-8}$. We implement a linear warm-up over 2000 steps, followed by a cosine annealing schedule driving the learning rate toward zero. Regularization is applied via a weight decay coefficient of 0.1, while gradients are globally clipped to a maximum norm of 1.0.

\section{Scaling Trends}
\label{section:scaling}

We offer a comprehensive overview of scaling behaviors exhibited by byte-level models, with the objective of guiding subsequent scaling efforts for {\textsc BLT} architectures. Our analysis seeks to overcome previous limitations in the study of byte-level models through several key approaches: (a) assessing trends within compute-optimal training protocols, (b) training comparable 8B-parameter models using substantial datasets (extending to 1T tokens or 4T bytes) and examining their performance on downstream evaluation tasks, and (c) evaluating scaling dynamics under conditions where inference cost is kept constant. In subsequent sections, we will further explore particular strengths associated with modeling byte-sequences.

\subsection{Parameter Matched Compute Optimal Scaling Trends}
\label{section:parameter-matched-scaling}
Utilizing the Llama 2 dataset, we conduct training of both \textit{compute-optimal} bpe and {\textsc BLT} models at four distinct scales, with model sizes spanning from 1B to 8B parameters. Subsequently, we visualize the relationship between the number of training flops and language modeling performance on a representative portion of the training data composition. The bpe models follow the optimal model size-to-training data ratio established by Llama~3~\citep{dubey2024llama}, adhering to the \textit{compute-optimal} regime intended to yield maximal performance for a fixed training compute budget~\citep{hoffmann2022training}. This serves as a reliable reference point for model comparison. For each bpe model, we also train a {\textsc BLT} counterpart on the identical dataset, employing a Latent Transformer whose architectural characteristics and parameter count mirror those of the corresponding bpe Transformer.

As shown in~\autoref{fig:scaling-compute-opt} (right), {\textsc BLT} models consistently equal or surpass the performance of their bpe equivalents as both model capacity and flops increase. To our knowledge, {\textsc BLT} represents the first byte-level Transformer architecture to reach scaling patterns comparable to BPE-based models when operating in compute-optimal settings. This observation supports our hypothesis that the parameter-to-training compute ratio optimal for bpe models is similarly effective for {\textsc BLT}, or at a minimum, does not diverge significantly.

To effectively align with bpe scaling patterns, it is essential to incorporate both enhancements in architecture and dynamic patching. In ~\autoref{fig:scaling-compute-opt} (left), we present a comparison between models utilizing space-patching and Llama 3. The performance of SpaceByte \citep{slagle2024spacebyte} is estimated by implementing {\textsc BLT} space-patching, omitting n-gram embeddings and cross-attention. While SpaceByte offers gains over Megabyte, its capabilities still fall significantly short of those of Llama 3. Meanwhile, \autoref{fig:scaling-compute-opt} (right) displays the cumulative benefits resulting from architectural updates in combination with dynamic patching. When scaled, {\textsc BLT} models achieve parity with leading tokenizer-based models like Llama 3.

Furthermore, our results indicate that for models relying on tokenizers, those utilizing the Llama-3 tokenizer achieve superior performance compared to counterparts trained with the Llama-2 tokenizer when provided with identical training data.

In summary, the {\textsc BLT} architecture exhibits intermediate performance relative to Llama 2 and Llama 3 when larger patch sizes are employed. While the Llama 2 and Llama 3 bpe tokenizers produce tokens averaging 3.7 and 4.4 bytes, respectively, {\textsc BLT} matches these models’ scaling properties using mean patch sizes of 6 or even 8 bytes. Since the inference flop count is inversely related to the patch size, adopting a patch size of 8 bytes allows for nearly a 50\% reduction in inference flops. Additionally, models that utilize greater patch sizes show improved performance as model capacity and dataset size are increased. For example, while {\textsc BLT} with an 8-byte patch begins with lower initial performance than the bpe-based Llama 2 at the 1B parameter scale, it surpasses bpe’s results at the 7B scale. This pattern indicates that larger patch sizes may yield superior performance as both model size and available computational resources continue to expand, suggesting that further increases in patch size might be practical and beneficial in the future.

\subsection{Beyond Compute Optimal Task Evaluations}
\label{section:evals}
To further investigate scaling characteristics, we extend the training of an 8B {\textsc BLT} model past the compute-optimal threshold using the {\textsc BLT}-1T dataset, which is both larger and of higher quality. We then evaluate the resulting model on a range of established classification and generative benchmarks. For this evaluation, we select tasks designed to test common sense reasoning, general world knowledge, and code generation proficiency:
\paragraph{Classification tasks} consist of ARC-Easy (0-shot)~\citep{Eval_ARC}, ARC-Challenge (0-shot)~\citep{Eval_ARC}, HellaSwag (0-shot)~\citep{Eval_hellaswag}, PIQA (0-shot)~\citep{Eval_piqa}, and MMLU (5-shot)~\citep{hendrycks2020measuring}. To evaluate performance, we use a prompt-scoring technique that computes the probability assigned to each candidate answer, and present mean accuracy across tasks.
\paragraph{Coding related generation tasks:}  We evaluate the capacity of LLMs to produce Python code by providing pass@1 metrics on MBPP (3-shot)~\citep{Eval_MBPP} and HumanEval (0-shot)~\citep{Eval_HumanEval}.

\autoref{tab:evals} presents a comparison among three models, all trained on the {\textsc BLT}-1T dataset: a bpe Llama 3 tokenizer-based model,\footnote{The Llama 3 tokenizer, featuring a 128k vocabulary, is selected for its superior performance over the 32k vocabulary of Llama 2.} and two distinct {\textsc BLT} model variants. Specifically, one uses a space-patching approach ({\textsc BLT}-Space), while the other adopts an entropy-based patching method ({\textsc BLT}-Entropy), applying an approximate monotonicity constraint and resetting the context upon encountering new lines in the entropy-based model (see~\autoref{section:entropy-context} for details). Each model is allotted the same computation (flop) budget during training. Notably, for the {\textsc BLT}-Entropy model, we implement an adjustment at inference: lowering the entropy threshold from 0.6 to 0.1. This change yields enhanced task performance, although it results in an increased number of inference steps.

The {\textsc BLT}-Entropy model surpasses the Llama 3 model in performance on 4 of the 7 evaluated tasks, despite both models being trained with an equivalent byte count. This enhancement can likely be attributed to (1) more efficient utilization of training compute achieved through dynamic patching, and (2) the explicit modeling of information at the byte level rather than the token level.

Conversely, {\textsc BLT}-Space generally lags behind the Llama 3 tokenizer across most tasks, except for one; however, it notably decreases inference flops due to its greater mean patch size of 6 bytes. In contrast, the models employing bpe and entropy-patching methodologies achieve a similar mean patch size, around 4.5 bytes, when evaluated on the training data mixture. Given an equal training allocation, the larger patch size of {\textsc BLT} allows it to process 30\% more data compared to the other two approaches, potentially moving {\textsc BLT} farther from compute-optimal efficiency.

\subsection{Patches Scale Better Than Tokens} The {\textsc BLT} architecture enables simultaneous scaling of both the model and patch dimensions, all while preserving a fixed computational budget for training and inference, as well as using a constant amount of training data. The capacity to arbitrarily expand patch sizes distinguishes patch-based models from traditional token-based approaches, allowing them to overcome efficiency limitations inherent to fixed-vocabulary tokenizers, as outlined in Section~\ref{section:bpe-patch}. By utilizing larger patches, computational demands are reduced, thereby freeing resources that can be allocated to increasing the capacity of the global latent transformer, which, as a result, needs to be executed fewer times.

We perform a fixed inference scaling analysis to evaluate whether increasing model size—paired with a reduced number of steps on larger patches—yields superior performance compared to deploying smaller models that execute more steps. Using baseline models with 400 million and 3.6 billion parameters and employing the Llama 2 tokenizer, we identify models with comparable compute requirements utilizing the Llama 3 tokenizer, as well as {\textsc BLT}-Entropy architectures featuring average patch sizes of 6 and 8 bytes on the training datamix (refer to~\autoref{tab:fixed-inf-params} for specifics regarding model configurations). For models with a patch size of 8, we apply 3 encoder layers rather than 1. Each configuration is trained using several distinct training flop allocations.

As illustrated in \autoref{fig:fixed_inference_scaling}, {\textsc BLT} models demonstrate superior scaling behavior compared to tokenization-based models across both categories of inference flops. While BPE models exhibit stronger performance under limited training budgets, they are soon outperformed by {\textsc BLT} as the total compute approaches and moves beyond the compute-optimal point. In real-world scenarios, allocating additional resources to pretraining can result in models that offer higher performance for the same fixed inference cost. This is exemplified by the 8B model category, such as Llama 3.1, which underwent pretraining on substantially more data—by nearly two orders of magnitude—than is considered compute-optimal for its model size.

The point at which {\textsc BLT} begins to outperform token-based models shifts somewhat nearer to the compute-optimal threshold as models of greater flop capacity are considered, moving from a margin of 3x to approximately 2.5x above the compute-optimal allocation. Additionally, the model utilizing a patch size of 8 demonstrates a more pronounced scaling advantage in this higher flop regime, surpassing the performance of alternative models at an earlier stage. As elaborated in Section~\ref{section:parameter-matched-scaling}, models with increased patch sizes tend to show performance approaching that of BPE models when scaled to larger model sizes. We partly explain this observation by noting that the byte-level Encoder and Decoder modules constitute a diminishing fraction of the total computational cost, as these components scale up at a slower rate compared to the Latent Transformer. For instance, when increasing the overall parameter count twentyfold—from 400M to 8B—only about twice as many local parameters are attributed to the {\textsc BLT} subcomponents. This distinction is noteworthy, as larger patch sizes impact the computational requirements of the patch Latent Transformer but not those of the byte-level modules. Indeed, it is for this reason that the model {\textsc BLT}-Entropy with ps=8 experienced an increase from 1.6x to 1.7x relative to the Llama 2 model size as model scale increased.

In conclusion, our investigation into patch-length scaling reveals that the {\textsc BLT} architecture benefits from concurrent enlargement of both the patch size and the model size, yielding superior scaling behavior. Notably, these improvements persist and may even be enhanced as model sizes become larger.

\section{Byte Modeling Improves Robustness} Additionally, we evaluate the robustness of {\textsc BLT} relative to models based on tokens without intrinsic byte-level knowledge, and introduce a strategy to convert pretrained token-centric models into a byte-aware format.
\label{sec:robustness}

\subsection{Character-Level Tasks}

Initial interest in byte-level models stemmed from their inherent resilience to noise present at the byte level in input data, as well as their capability to recognize token composition—an area where contemporary tokenizer-based architectures often face limitations. To assess these characteristics, we conduct supplementary experiments using benchmarks specifically designed to test both resilience to input disturbances and understanding of character-level information, covering English, multiple languages, digits, and phonemes. The outcomes of these experiments are summarized in Table~\ref{tab:char_tasks}.

\paragraph{Noisy Data} To assess the robustness of {\textsc BLT} relative to tokenizer-based models, we construct perturbed versions of the classification benchmarks detailed in Section~\ref{section:evals}. We utilize five unique methods at the character level to introduce noise into the input text: (a) \textit{AntSpeak}: Here, the text is converted such that every character is uppercase and characters are separated by spaces. (b) \textit{Drop}: This approach randomly eliminates 10\% of all characters present in the text. (c) \textit{RandomCase}: Throughout the input, 50\% of characters are randomly assigned uppercase and 50\% lowercase casing. (d) \textit{Repeat}: In this variant, 20\% of characters are repeated between two to four times. (e) \textit{UpperCase}: This simply turns all letters to uppercase. Each noising procedure is individually applied to the prompt, the completion, or to both elements, and performance is averaged across these cases during evaluation. Table \ref{tab:char_tasks} provides a comparison of results on the noised HellaSwag~\citep{Eval_hellaswag} dataset, revealing that {\textsc BLT} consistently displays stronger robustness compared to tokenizer-dependent models—exceeding by an average of 8 points models trained on equivalent data, and outperforming the Llama 3.1 model that benefits from substantially more training data.

\paragraph{Phonology - Grapheme-to-Phoneme (G2P)} We evaluate {\textsc BLT}'s proficiency in converting sequences of graphemes (the written symbols of a word) into their corresponding phoneme transcriptions, which represent pronunciation. Table \ref{tab:char_tasks} displays the 5-shot G2P task results using Phonology Bench~\citep{suvarna2024phonologybench}, demonstrating that {\textsc BLT} achieves superior performance compared to the Llama 3 1T tokenizer-based baseline on this evaluation.

\paragraph{CUTE}
To evaluate the model's capability for character-level reasoning, we test {\textsc BLT} using the CUTE benchmark~\citep{edman2024cute}, which contains a set of tasks grouped into three primary categories: comprehension of composition, recognition of orthographic similarity, and sequence manipulation skills. This benchmark is notably challenging for most models based on tokenization, which tend to know the orthographic makeup of their tokens but encounter difficulties in leveraging this information for manipulating text sequences. As seen in Table~\ref{tab:char_tasks}, {\textsc BLT}-Entropy surpasses both BPE Llama 3 variants by over 25 points on this suite of tasks. Notably, our model excels in tasks centered on character manipulation, attaining 99.9\% accuracy on both spelling categories. These substantial gains, achieved even though {\textsc BLT} was trained on sixteen times less data than Llama 3.1, suggest that BPE models face inherent challenges acquiring robust character-level representations. Figure \ref{fig:cute_examples} presents examples where Llama 3's tokenizer-based model falters, whereas {\textsc BLT} succeeds. Of the task suite, only word deletion and word insertion show superior performance from BPE-based models. While byte-level approaches like {\textsc BLT} may find explicit word manipulation less straightforward, the performance gap remains relatively small, and composing word-level understanding from characters may be more tractable than inferring fine-grained character patterns from word tokens. Our evaluation maintains consistent methodology and uses original Huggingface prompts for all experiments. With targeted prompt engineering, BPE-based models could potentially realize additional gains.

\paragraph{Low Resource Machine Translation} We assess the performance of {\textsc BLT} in both directions of translation across six widely spoken language families and twenty-one lower-resource languages, each utilizing different scripts, using the FLORES-101 benchmark~\citep{goyal2022flores}. SentencePiece BLEU results are provided in Table~\ref{tab:flores}. The findings reveal that {\textsc BLT} surpasses a model employing the Llama 3 tokenizer, yielding an overall BLEU improvement of 2 points for translations into English, and a 0.5 point increase for translations originating from English. For commonly studied language pairs, {\textsc BLT} delivers performance on par with, or marginally superior to, Llama 3. Notably, across a variety of lower-resource language families, {\textsc BLT} registers clear gains over Llama 3, highlighting the strength of byte-level modeling in effectively capturing patterns within low-frequency byte sequences.

\subsection{Training {\textsc BLT} from Llama 3}
\label{sec:distillation}
We investigate an approach wherein {\textsc BLT} models benefit from pre-existing, pre-trained tokenizer-based models to enhance both training efficiency and convergence. This is accomplished by initializing the global transformer weights in {\textsc BLT} with those from a pre-trained Llama 3.1. Following this initialization, we update the global transformer’s parameters at a learning rate one-tenth that of the learning rate applied to the local encoder and decoder, following the optimal number of training steps identified for Llama 3. We present a comparative analysis with a standard {\textsc BLT} model in Table~\ref{tab:distill}. The results clearly show that initializing {\textsc BLT} with Llama 3.1’s weights yields substantial performance gains over both the original Llama 3 and conventional {\textsc BLT} baselines when matched for total computational cost (in flops). In addition, this version of {\textsc BLT} surpasses our {\textsc BLT}-Entropy variant (see Table~\ref{tab:evals}), even though the latter was trained on a much larger corpus (1T tokens). This indicates that the Llama 3.1-based initialization enables significant reductions in required training flops, while delivering superior MMLU task results.

This approach can alternatively be interpreted as converting models that rely on tokenizers into models that do not, thereby transforming a pre-trained LLaMA 3.1 into a {\textsc BLT} architecture. For a thorough evaluation, Table \ref{tab:distill} presents results for both the original LLaMA 3.1—trained on 15T tokens—and the {\textsc BLT} model adapted from LLaMA 3. Empirical results reveal that our model suffers a modest decrease in accuracy on MMLU and HumanEval, while the performance loss is more pronounced across additional tasks. These observations highlight the necessity for future efforts to more fully exploit the capabilities of the pre-trained backbone, particularly by refining data mixture strategies and adjusting other critical hyperparameters.

\section{Ablations and Discussion}
\label{section:ablations}
Here, we present a series of ablation studies that examine the rationale behind the architectural design of {\textsc BLT}, alongside the selection of patching methodologies and hyper-parameters for the {\textsc BLT} 8B-parameter model developed using the {\textsc BLT}-1T dataset.

\paragraph{Entropy Model Hyper-parameters} To investigate how the size of the entropy model and the context window length influence scaling, we conduct experiments with byte-level entropy transformer models whose parameter counts range from 1 million up to 100 million, and whose context window sizes vary between 64 and 512. We then chart the relationship between bits per byte (bpb) and training FLOP by plotting scaling law curves using our $400m$ and $1b$ {\textsc BLT} models, both trained on the Llama-2 dataset, as shown in \autoref{fig:entropy_ablation}. The results indicate that improving either the entropy model size or the context window enhances scaling behavior, although performance gains tend to taper off once the model exceeds 50 million parameters.

\paragraph{Types of Patching}

We conduct an ablation study of the four patching methods outlined in Section~\ref{section:patching}: (1) Strided Patching with strides set to 4 and 6; (2) Patching at whitespace positions; (3) Patching via the BPE Tokenizer aligned with the Llama 3 tokenizer; and (4) Entropy-based patching leveraging a compact byte-level llm.

Although dynamic patching shortens the effective sequence length, we standardize sequence length across all patching strategies to ensure a consistent context length. Consequently, every model is exposed to an equivalent quantity of bytes per sequence on average during both training and inference, eliminating potential confounds related to modeling longer contexts. The ablation results are illustrated in Figure~\ref{fig:scaling-compute-opt}. Notably, all alternative patching strategies surpass static patching, with space patching exhibiting performance nearly on par with dynamic entropy based patching.

In \autoref{tab:evals_ablation}, we provide results from benchmark assessments of {\textsc BLT} models, making comparisons among models utilizing tokenizers, space patching, and entropy-based patching techniques, each trained on the Llama 2 dataset for the optimal training duration~\citep{dubey2024llama}. Despite space patching being a more straightforward approach that avoids the computational overhead of executing an entropy model during training, our results demonstrate that the improvements associated with entropy-based patching, previously identified in scaling analyses~(Section~\ref{section:scaling}), also manifest in downstream benchmark evaluations.\footnote{Space patching outcomes are derived from previous experiments conducted without cross-attention; however, the observed patterns remain consistent when cross-attention is included.}

\paragraph{Cross-Attention}
As presented in~\autoref{tab:crossattn_ablations_1b}, we analyze the impact of inserting cross-attention at different locations within both the encoder and decoder components of {\textsc BLT}. Specifically for the encoder's cross-attention, we examine three approaches to initializing the queries: (1) assigning an identical learned embedding to all global states, (2) employing a hash embedding derived from the patch bytes, and (3) aggregating the encoder's hidden states corresponding to patch bytes at the specified encoder layer.

Our results indicate that incorporating cross-attention within the \textit{decoder} yields the strongest performance gains. In contrast, applying cross-attention in the encoder offers only marginal benefits, and these occur exclusively when queries are initialized through pooling. Furthermore, our analysis shows that cross-attention is especially advantageous for Common-Crawl data, with its effect becoming more pronounced as patch size increases.

\paragraph{n-gram Hash Embeddings} 

We conduct ablations using n-gram hash embedding vocabulary sizes of 0, 100K, 200K, and 400K, with corresponding results summarized in Table~\ref{tab:ngram_ablations_1b}. Our analysis shows that incorporating hash embeddings consistently benefits performance across all domains, with the most notable improvements observed on Wikipedia and Github (demonstrated by a 0.04 bpb improvement versus a 0.01 bpb gain at 8B after 15k steps). At the 8B scale, reducing the number of hashes from 500K to 300K leads to only a marginal decrease of approximately 0.001 bpb after 15k steps. These findings underscore the importance of hash embeddings for enabling {\textsc BLT} models to achieve performance comparable to tokenizer-based models, although performance improvements plateau beyond 300K hashes. Furthermore, the enhancements provided by hashes and cross-attention tend to be complementary, as each mechanism yields performance gains on separate datasets.

\paragraph{Local Model Hyperparameters} Table \ref{tab:local_layers} presents an ablation study on different configurations for the number of layers in the local encoder and decoder. Notably, when combined with hash n-gram embeddings, {\textsc BLT} achieves strong performance using a highly lightweight encoder—consisting of only a single layer—alongside a more substantial, multi-layered decoder.

\section{Related Work}
\label{section:related_work}

\textbf{Character-Level RNNs:} The task of Character Language Modeling has held significant interest since the initial development of neural approaches~\citep{sutskever2011generating,mikolov2012subword,graves2013generating}, in large part because such models can natively represent and generate out-of-vocabulary words without relying on explicit back-off strategies. \cite{kim2016character} introduced a model utilizing character-level processing only on the input, combining convolutional layers and highway networks with LSTM-based RNNs. Their architecture achieved parity with leading RNN-based language models in English, and surpassed them for languages with complex morphological structure—highlighting a notable advantage of character-level large language models (LLMs). In related work, \cite{kenter2018byte} demonstrated that byte-level LSTM architectures provide superior performance in machine comprehension tasks over word-based alternatives, particularly for Turkish and Russian, both morphologically rich languages. Expanding on these concepts, \cite{zhang2015character} leveraged convolutional networks at the character level for classification, achieving results that sometimes exceeded those of word-level systems. Hierarchical approaches were proposed by \citet{chung2019hierarchical}, who employed multi-level LSTMs with boundary detectors to uncover textual hierarchies, which led to further improvements in character-level language modeling. Distinct from attention-based methods, ByteNet~\cite{kalchbrenner2016neural} employed character-level convolutional networks for the task of machine translation.

\textbf{Character-Level Transformers:} The introduction of attention-based transformer models~\citep{vaswani2017attention} alongside subword tokenization techniques~\citep{sennrich-etal-2016-neural} greatly advanced the effectiveness of neural architectures on language modeling and established benchmarks. Despite these achievements, both word and sub-word representations introduce an inherent inductive bias regarding the model's level of linguistic abstraction. In order to integrate the strengths of transformer architectures with early encouraging results from character-level language modeling, \citet{al2019character} implemented particularly deep transformer networks and incorporated auxiliary objectives. Their transformer-based approaches managed to surpass previous character-level LSTM models in terms of performance. Nevertheless, there remained a considerable performance discrepancy relative to word-level large language models. Moreover, the findings from GPT-2~\citep{radfordgpt2} indicated that on extensive datasets such as the 1 billion word benchmark, models operating at the byte level lagged behind their word-level counterparts in competitiveness.

Although \cite{Choe2019BridgingTG} showed that transformer-based byte-level LLMs can exceed the performance of subword-level counterparts with similar parameter counts, these models entail substantially greater computational requirements and lengthier training times. In a related vein, \cite{el2020characterbert} introduced CharFormer, a BERT-style architecture that constructs word representations through convolutions applied to character embeddings, yielding improved results in the medical domain at the cost of significantly more computation. The CANINE model, developed by \cite{clark2022canine}, is a 150M parameter encoder-only system that directly handles character sequences. CANINE employs a deep transformer architecture akin to the global component of our approach and integrates local transformers with strided convolutional downsampling to reduce the length of character sequences. This model surpasses its token-level encoder-only analogue, mBERT, across multiple multilingual downstream benchmarks. ByT5~\citep{xue2022byt5} investigated fully byte-level encoder-decoder architectures, intentionally avoiding any form of patching. Although these models demonstrated increased noise robustness and rivaled tokenization-based systems with four times less data, the absence of patching led to significant computational overhead due to the necessity of executing attention over all bytes in an input. Representing inputs at the byte rather than subword granularity inherently increases sequence lengths, impeding the computational scalability of byte-level architectures. More recently, leveraging the fixed-size memory capabilities of the Mamba Architecture~\citep{gu2023mamba} over extended context windows, \cite{wang2024mambabyte} trained a byte-level Mamba model without patching, achieving superior performance over byte-level transformers in flop-controlled evaluations at the 350M parameter scale, as measured by bits-per-byte across several datasets.

\textbf{Patching-based approaches:} Leveraging patching can help mitigate the high computational cost associated with byte-level LLMs without necessarily compromising model performance. Numerous studies have shown promising outcomes with such techniques, particularly on smaller models and limited datasets. For instance, \cite{nawrot-etal-2022-hierarchical} implement static patch-based downsampling and upsampling strategies, introducing the hourglass transformer architecture, which achieves superior results compared to other byte-level benchmarks at the 150M parameter level. Building on this, \cite{nawrot-etal-2023-efficient} introduce enhancements via dynamic patching mechanisms, such as an end-to-end learnable boundary-predictor, a boundary-predictor guided by specific tokenizers, and an entropy-driven patching scheme akin to {\textsc BLT}. They demonstrate that these innovations can surpass standard transformers of the same period in language modeling tasks, specifically with models at the 40M parameter scale on approximately 400M tokens. Separately, \cite{lester2024training} examine training with sequences that are compressed through arithmetic coding, thus achieving higher compression rates than BPE. By employing an equal-info windows approach, they exceed byte-level baseline performance on language modeling tasks, although results still lag behind those of subword-based baselines.

Our research is heavily influenced by and most directly related to MegaByte~\citep{yu2023megabyte}, a decoder-only causal language model that implements a fixed, static approach to segmenting and joining byte representations into patches. MegaByte also employs a dedicated local model within the decoder to revert patches back to byte sequences. The creators of MegaByte demonstrate its ability to achieve performance on par with tokenizer-based architectures at the 1B parameter scale over a dataset comprising approximately 400B bytes. In our investigations, we consistently benchmark MegaByte and observe that its reliance on static patching results in suboptimal performance compared to the latest compute-efficient, tokenizer-based models under comparable FLOP budgets. Furthermore, our work illustrates how {\textsc BLT} can effectively close this performance gap. In a similar vein, \cite{slagle2024spacebyte} independently arrive at analogous conclusions concerning MegaByte, proposing adaptations that involve patching at white spaces and other space-like byte boundaries and introducing a local encoder. Their findings indicate that such modifications can provide performance gains over tokenizer-based transformer models within certain domains—such as Github and arXiv—at the 1B parameter scale when FLOPS are held constant. We include empirical results involving this approach and argue that substantial advances in architecture are still required to enable byte-level models to scale and reach parity with leading token-based models, exemplified by systems like Llama 3.

\section{Limitations and Future Work}

For this study, we adopt the number of training steps found optimal for Llama~3~\citep{dubey2024llama} when selecting architectural configurations. Nonetheless, it is important to note that these scaling laws were originally derived for BPE-level transformer architectures, which may result in less-than-ideal relationships between data and parameter sizes when applied to {\textsc BLT}. The task of deriving scaling laws specific to {\textsc BLT}—potentially revealing even more advantageous scaling characteristics for this architecture—remains as an avenue for future investigation. Furthermore, as many of our experiments have been limited to models with up to 1B parameters, there remains the possibility that scaling up to 8B parameters or higher could shift the optimal architectural choices, potentially yielding superior performance at these larger model sizes.

Current transformer codebases and libraries are heavily optimized for models that rely on tokenizers and follow the standard transformer framework. Although we conduct experiments that match theoretical flop counts and leverage some efficient solutions—like FlexAttention—to accommodate architectural modifications beyond the classic transformer, our implementations may still fall short of tokenizer-based models in actual runtime performance and could improve with additional optimizations.

Although {\textsc BLT} currently employs a distinct, pre-trained entropy model for patching, an intriguing avenue for future research lies in training the patching component in an end-to-end manner. Section \ref{sec:distillation} discusses preliminary studies that provide promising signs of success in ``byte-ifying'' models such as Llama 3, which are tokenizer-based and have been trained on datasets exceeding 10T tokens; these experiments initialize and freeze the global transformer using the original weights. Continued exploration along these lines may yield strategies that preserve the advantages of bytefying, while potentially surpassing the capabilities of the original tokenizer-dependent architectures—without requiring the models to be retrained from the ground up.

\section{Conclusion}
\label{section:conclusion}

This work introduces the Byte Latent Transformer (\textbf{BLT}), an innovative architecture that moves beyond the standard reliance on fixed-vocabulary tokenization typical in large language models. {\textsc BLT} incorporates a trainable, adaptive approach for aggregating bytes into patches, which enables the architecture to dynamically allocate computation depending on the intricacy of input data. This design leads to marked gains in both computational efficiency and model robustness. Comprehensive scaling experiments reveal that {\textsc BLT} models are able to achieve parity with token-based models such as Llama 3 at comparable parameter counts, up to 8B parameters and processing 4T bytes, while offering the capability to reduce inference computation (flops) by up to 50\% with only minimal sacrifices in benchmark performance. Additionally, {\textsc BLT} opens a novel path for scaling models: increasing both the model's capacity and the size of its input patches concurrently within a fixed inference cost. This approach offers distinct advantages under computational constraints frequently observed in real-world applications. Beyond simply processing at the byte level, {\textsc BLT} enhances model resilience to input noise and improves its grasp of fine-grained linguistic structures, substantially benefiting long-tail data scenarios. Altogether, these findings position {\textsc BLT} as a compelling alternative to established tokenization-based techniques, delivering a flexible, scalable, and robust solution for modern language modeling.

\end{document}